\section{Introduction}\label{sec:intro}

% TODO(BH): draft. Raw material: JournalPaper/positioning_memo.md §2
% (motivation stack). Suggested skeleton (logic only, prose is yours):
%
% 1. Phenomenon: DeFi exploit losses in the billions annually; smart-
%    contract risk fails classical insurability criteria (Bekemeier 2023);
%    dedicated cover protocols remain marginal; the recognized bottleneck
%    is the yield conundrum (Zhou 2026).
% 2. The tension: insurance under ADVERSARIAL risk. Losses are chosen by
%    a profit-motivated attacker who responds to the contract; evidence
%    that attackers condition on insurance (Meurs et al. 2022: insured
%    victims pay 2.7x larger ransoms; Li & Liao 2025: insurance raises
%    attacker payoffs).
% 3. The Budish gap: incentive-based security requires attack cost
%    scaling ~linearly with value secured (Budish 2025); our structural
%    estimate of the application-layer technology is eta ~ 0.2. The gap
%    between required and actual scaling is the market for insurance.
% 4. This paper: a coverage-pool mechanism whose two DeFi-specific
%    instruments -- negligence-contingent collateral forfeiture and an
%    above-market pool yield -- jointly neutralize attacker-amplified
%    moral hazard. Preview Prop. D (decoupling), Prop. C (carry as
%    security subsidy), and the calibrated counterfactuals (forfeiture
%    off: hacks 2.3 -> 10.8/yr).
% 5. Why blockchain: smart contracts make forfeiture credible and
%    automatic (Corollary D1).

\emph{[To be drafted.]}
